\subsection{Attacks on the Implementation of Cryptographic Algorithms}\label{sec:security_of_cryptographic_algorithms:attacks}

We now present those 3 categories of attacks (already mentioned in \Cref{sec:cryptographic_security_overview:attacks}) in which the way an algorithm is implemented is relevant, i.e., side-channel attacks (\Cref{sec:security_of_cryptographic_algorithms:attacks:sca}), fault injection attacks (\Cref{sec:security_of_cryptographic_algorithms:attacks:fia}), and data remanence attacks (\Cref{sec:security_of_cryptographic_algorithms:attacks:dra}).

\remarkBlock{Classification of side-channel attacks}{Depending on the criteria used to distinguish different categories of attacks, fault injection attacks and data remanence attacks are sometimes seen as a specific instance of side-channel attacks (e.g., see \cite{piessens_side_2024,spreitzer_systematic_2018,munoz_survey_2023}).}

\subsubsection{Side-Channel Attacks}\label{sec:security_of_cryptographic_algorithms:attacks:sca}

At a high level, \textbf{side-channel attacks} measure side effects --- generally referred to as \textbf{leakage} --- of cryptographic computations to infer (portions of) secret data such as seeds, private keys, and internal states. The term side-channel attack is an \textbf{umbrella term} that encompasses several --- and possibly fundamentally different --- approaches for attacking the implementation of algorithms by gathering (separately or collectively) various kinds of leakages: timing, power consumption, radio and electromagnetic leaks, and even acoustic emissions, to mention a few. Nonetheless, all side-channel attacks are generally carried out by:

\begin{enumerate}
	\item identifying secret data-dependent cryptographic computations;
	\item recording traces caused by such computations;
	\item analyzing such traces to recover secret data.
\end{enumerate}

While side-channel attacks are typically passive attacks, some may be invasive in the sense that they require accessing cryptographic modules (\Cref{sec:cryptographic_modules}) and operating on them (e.g., insert a probe on a wire to read electric signals) \cite{noauthor_power_2007}.

\remarkBlock{\gls{snr}}{A central concept in the effectiveness of side-channel attacks is the \textbf{\gls{snr}}, which quantifies the relationship between useful information (\textbf{signal}) and irrelevant data or interference (\textbf{noise}) within the monitored leakage. In other words, the \gls{snr} reflects how distinguishable the useful information is from the irrelevant data or the interference: a high \gls{snr} means the signal is prominent and therefore easy to isolate, while a low \gls{snr} means the noise is prominent, posing greater difficulty in isolating the signal. Consequently, side-channel attacks aim to maximize the \gls{snr}, while mitigations (\Cref{sec:security_of_cryptographic_algorithms:best_practices}) aim to either eliminate or reduce the signal (e.g., by reducing sources of signal or eliminating the relationship between the signal and the secret data, respectively) or increase the noise.}

\paragraph*{\faBolt~Power Analysis Attacks.} This type of side-channel attack relies on measuring the \textbf{power consumption} of a cryptographic module during the execution of cryptographic computations. In fact, depending on how algorithms are implemented, secret data (e.g., private keys) may influence power consumption, therefore creating correlations that an attacker may analyze. Intuitively, attackers must have access to the cryptographic module and often be able to control its execution to conduct power analysis attacks \cite{messerges_investigations_1999} --- hence, attackers must be local (\Cref{sec:cryptographic_security_overview:threat_model}). At a high level, power analysis attacks may be \cite{goos_differential_1999,goos_towards_1999} \textit{simple} or \textit{differential}: in the former case, single-trace local attackers try to directly (and often visually) correlate power consumption with cryptographic computations to reveal secret data, while in the latter case multi-trace local attackers rely on statistical methods (e.g., average) to correlate small variations in power consumption with secret data while collecting multiple measurements for different inputs. Simple power analysis attacks usually require a high \gls{snr} and are effective when secret data-dependent behavior in power consumption is apparent, while differential power analysis attacks --- by using statistical methods (as well as signal processing and error correction) --- are effective even with low \gls{snr}.

\remarkBlock{Latest trends for power analysis attacks}{Starting from the seminal work in \cite{goos_differential_1999}, several power analysis attacks were discussed in the literature such as the Collide+Power attack \cite{kogler_collide_2023},\footnote{\url{https://collidepower.com/}}. Recently, many power analysis attacks are combining power analysis with machine learning and artificial intelligence \cite{liran_power_2014} and targeting post-quantum algorithms \cite{kamucheka_power_2021}}.
%TODO (\Cref{future_sec:post_quantum_cryptography}).}

\paragraph*{\faBroadcastTower~Radio and Electromagnetic Analysis Attacks.} This type of side-channel attack relies on capturing radio or electromagnetic emissions from hardware to infer currently ongoing cryptographic computations and secret data within cryptographic modules. Radio and electromagnetic analysis attacks are strictly related to power analysis attacks, especially because \textbf{radio and electromagnetic emissions are a direct consequence of power consumption}. It follows that radio and electromagnetic analysis attacks expect an analogous threat model (i.e., require local attackers), use techniques equivalent to those of power analysis attacks, and are similarly divided between simple and differential. However, radio and electromagnetic analysis attacks are more complex and difficult to carry out, having achieved lower maturity \cite{sayakkara_survey_2019}. Conversely, sampling and probing for radio and electromagnetic analysis attacks is easier than for power analysis attacks, as the former requires neither physical access nor depackaging/decapsulation (i.e., exposing the silicon of chips).

\paragraph*{\faMicrophone~Acoustic Analysis Attacks.} This type of side-channel attack relies on measuring \textbf{ultrasonic sounds and vibrations} emitted from hardware to then infer secret data. Such ultrasonic sounds and vibrations are usually not human-audible and come from internal electronic components (e.g., capacitors and inductors) \cite{genkin_acoustic_2017}. Similarly to power, radio, and electromagnetic analysis attacks, acoustic analysis attacks may be simple and differential. 

\remarkBlock{Threat model for acoustic analysis attacks}{Acoustic analysis attacks may not always require attackers to be local when cryptographic modules are operating near attacker-controlled microphones; this is often the case when considering cryptographic modules within laptops.}

\paragraph*{\faThermometerHalf~Thermal Imaging Attacks.} This type of side-channel attack relies on the observation of \textbf{residual heat and heat patterns} on hardware --- usually through thermal imaging cameras capturing infrared images --- to then infer secret data. Intuitively, thermal imaging attacks, regardless of whether they are simple or differential, require local attackers. Like other side-channel attacks (i.e., acoustic, radio, and electromagnetic) whose traces are linearly correlated with power consumption, thermal imaging attacks are complex and difficult to carry out, especially if considering the limitations posed by the physical properties of thermal conductivity and capacitance \cite{francillon_temperature_2014}. For this reason, thermal imaging attacks are often carried out in combination with other kinds of side-channel attacks \cite{aljuffri_applying_2021}.

\paragraph*{\faCameraRetro~Optical Analysis Attacks.} This type of side-channel attack relies on the phenomena of \textbf{photonic emission} from
transistors when executing cryptographic computations \cite{ferrigno_when_2008,hutchison_simple_2012} to then infer secret data. Optical analysis attacks require local attackers and may be simple or differential. In terms of equipment, an optical analysis attack is comparable in price to a power analysis attack \cite{carmon_photonic_2017}.

\paragraph*{\faStopwatch~Timing Attacks.} This type of side-channel attack relies on monitoring \textbf{variations in the execution time} of cryptographic computations to infer secret data \cite{kocher_timing_1996}. Timing attacks are among the most studies, effective, and impactful side-channel attacks. Even worse, timing attacks --- whether simple or differential --- can be conducted either locally and even remotely over the internet \cite{brumley_remote_2003,crosby_opportunities_2009}. 

\exampleBlock{A famous example of a timing attack}{In 1996, Paul Kocher demonstrated how to recover private \gls{rsa} keys (\Cref{sec:asymmetric_cryptography:ciphers:rsa}) via a timing attack --- see \cite{kocher_timing_1996}}

There are multiple reasons why cryptographic computations --- when implemented --- may take more or less time depending on secret data:

\begin{itemize}
	\item \textit{generic operations}: some (especially mathematical) operations may execute faster or slower according to the input data and their implementation, e.g., exponentiations for the \gls{dl} problem (\Cref{sec:trapdoor_functions:discrete_logarithm}) implemented with the square-and-multiply algorithm and fast-failing comparison functions like \texttt{memcmp()}\footnote{\url{https://en.cppreference.com/w/c/string/byte/memcmp}} in the C programming language;
	
	\item \textit{conditional branches and jumps}: making the execution (of a portion of the implementation) of a cryptographic algorithm dependent on secret data is one of the most common timing attack vulnerabilities. Moreover, to provide more concurrency when extra resources are available, modern processors often speculatively execute past branches and jumps by assuming the corresponding condition will be true. However, \textbf{speculative execution} either prevents (by performing operations in advance) or adds delays (by reverting changes) according to whether conditions are actually true. Especially when conditions are tied to secret data, speculative execution may result in timing and memory attacks (since processors may speculatively load secret data into memory before checking the process' permissions);
	
	\exampleBlock{Real-world speculative execution timing attacks}{Two famous attacks exploiting speculative execution are Spectre \cite{kocher_spectre_2019} and Meltdown \cite{lipp_meltdown_2020}. Note that, while being mainly based on speculative execution, Spectre and Meltdown also leverage leakage from cache behavior.}
	
	\item \textit{memory access patterns}: a large sub-category of timing attacks is based on memory access patterns and, in particular, \textbf{cache behavior}. At a high level, cache timing attacks rely on subtle timing variations in cache access patterns (e.g., cache hits and misses) which may ultimately depend on secret data. Intuitively, the way a cryptographic algorithm is implemented may make it more or less susceptible to cache timing attacks. 
	\exampleBlock{Real-world memory access patterns timing attacks}{\gls{aes} (\Cref{sec:symmetric_cryptography:ciphers:aes}) often uses lookup tables for efficiency. However, memory access patterns for these tables depend on the secret key, hence they may leak secret data as presented in \cite{tromer_efficient_2010,ashokkumar_highly_2016,stavrou_wait_2014}.}
	There exist many types of cache timing attacks, \textbf{all of which require an active attacker}; the main 4 are:
	
	\begin{itemize}
		
		\item \texttt{Evict+Time} \cite{hutchison_cache_2006} (and then in \cite{tromer_efficient_2010}): the attacker ensures specific cache lines are evicted (i.e., removed) from the cache. Then, when a cryptographic computation involving secret data is executed, the attacker measures the time taken by the execution of the cryptographic computation and determines --- according to the access time --- whether the required secret data (or instructions) were not in the cache. \texttt{Evict+Time} is often used when cache access patterns depend on secret data;
		
		\item \texttt{Prime+Probe} \cite{hutchison_cache_2006} (and then in \cite{tromer_efficient_2010}): the attacker fills (i.e., primes) the cache with their own data, ensuring the cache lines are in a known state. Therefore, when a cryptographic computation involving secret data is executed, it is likely that some of the attacker's data get evicted. The attacker can then measure the time it takes to access their own data, determining --- according to the access time --- what data was evicted. \texttt{Prime+Probe} is often used on shared caches in multi-core or multi-tenant environments (e.g., cloud computing);
		
		\item \texttt{Flush+Reload} \cite{yarom_flush_2014,stavrou_wait_2014}: the attacker flushes a specific cache line (i.e., removes it from the cache) using instructions like \texttt{clflush}. When a cryptographic computation involving secret data is executed, the attacker accesses that specific cache line again, determining --- according to the access time --- whether that cache line was evicted. We note that \texttt{Flush+Reload} is more precise than \texttt{Prime+Probe} and it is often used when specific memory accesses correlate with secret data-dependent computations;
		
		\item \texttt{Flush+Flush} \cite{gruss_flushflush_2016}: a more complex yet stealthier variation of the \texttt{Flush+Reload} attack where the attacker flushes the cache once, the cryptographic computation involving secret data is executed, and then the attacker flushes the cache again while calculating the time taken to perform the second flush --- the rationale being that the time spent in flushing depends on what cache lines has been loaded by the cryptographic computation.
				
	\end{itemize}
\end{itemize}

As a final remark, note that timing attacks can be introduced unintentionally with compiler optimizations and \gls{isa} hardware characteristics --- see \Cref{sec:programming_languages_and_software} for more details.

\readBlock{More practical information on side-channel attacks}{Soatok's blog post ``Soatok's Guide to Side-Channel Attacks'' and Jean-Philippe Aumasson's GitHub post ``Cryptocoding'' --- available at the links \url{https://soatok.blog/2020/08/27/soatoks-guide-to-side-channel-attacks/} and \url{https://github.com/veorq/cryptocoding} --- provide a good excursus on side-channel attacks with code examples.}

\subsubsection{Fault Injection Attacks}\label{sec:security_of_cryptographic_algorithms:attacks:fia}

At a high level, \textbf{fault injection attacks} --- firstly introduced in \cite{biham_differential_1997} --- aim to intentionally introduce hardware faults during the execution of cryptographic computations, usually to infer (portions of) secret data but also to bypass security checks or control flows (e.g., by flipping bits, skipping or repeating operations), cause incorrect data fetch (e.g., from the cache), and understand how a cryptographic module behaves under induced faults \cite{barenghi_fault_2012}. Similarly to side-channel attacks (\Cref{sec:security_of_cryptographic_algorithms:attacks:sca}), the term fault injection attack denotes a \textbf{category of attacks} encompassing different approaches for introducing faults, e.g, glitching attacks \cite{dutertre_review_2011}, laser and optical fault injection \cite{selmke_attack_2016}, electromagnetic pulses, and temperature variations \cite{chiu_sok_2023}. Fault injection attacks are carried out by local and active attackers and are obviously invasive (except for some software-based fault injection attacks), and --- similarly to many side-channel attacks --- can be either simple or differential. Although fault injection attacks do not always require direct contact to introduce faults, depackaging/decapsulation is sometimes needed. Finally, fault injection attacks commonly target hardware-based cryptographic modules like \gls{se} and \gls{tpm} (and similar embedded systems) by affecting either processors (i.e., computations) or memory (i.e., data).

\paragraph*{\faSignature~Glitching Attacks.} This type of fault injection attack relies on introducing faults as either \textbf{power or voltage glitches} \cite{bozzato_shaping_2019} --- that is, sudden and short-lived interruptions or fluctuations in the power supply of a circuit --- or as \textbf{clock glitches} --- that is, irregularities in clock signal resulting in the speed up or delay of computations --- to cause computational errors or timing issues \cite{endo_chip_2011}.

\readBlock{More information on glitching attacks}{The NCCgroup's series of blogposts ``An Introduction to Fault Injection'' --- available at the link \url{https://www.nccgroup.com/us/research-blog/an-introduction-to-fault-injection-part-13/} --- provides further details on fault injection attacks based on glitching.}

\paragraph*{\faBullseye~Laser and Optical Fault Injection.} This type of fault injection attack relies on directing high-precision and timely laser \cite{skorobogatov_optical_2003} and optical \cite{schmidt_optical_2007} beams on specific locations of a circuit to introduce faults through \textbf{localized heating or photoelectric effects}.

\paragraph*{\faMagnet~Electromagnetic Pulses.} This type of fault injection attack relies on emitting electromagnetic pulses (or creating electromagnetic fields) on specific locations of a circuit to \textbf{induce voltages or currents} and introduce faults. Electromagnetic pulses are usually less precise (and potentially more damaging) than laser and optical beams \cite{chiu_sok_2023}.

\paragraph*{\faFire~Temperature Variations.} This type of fault injection attack relies on altering the temperature (i.e., through extreme cooling or heating) of specific locations of a circuit to \textbf{alter electrical characteristics} and introduce faults \cite{francillon_temperature_2014}. The rationale behind temperature variation attacks is that electronic components --- like transistors and diodes --- may slow down, speed up, or change their behavior when subject to significant temperature variations.

\paragraph*{\faLaptopCode~Software-Based Fault Injection.} This type of fault injection attack relies on \textbf{exploiting software and firmware vulnerabilities and bugs to introduce faults remotely}. Software-based fault injection attacks are particularly worrisome as they do not require local attackers, hence can even target cloud-hosted servers and virtual machines \cite{murdock_plundervolt_2020,zhang_cachewarp_2024}. The most common software-based fault injection attack is \textbf{Rowhammer} \cite{koksal_jolt_2002,kim_flipping_2014}, which aims to flip particular bits (belonging to a process performing cryptographic computations) in DRAM memory by repeatedly accessing nearby bits (belonging to a process controlled by the attacker) in adjacent rows of memory cells. The Rowhammer attack is feasible due to the fact that memory cells are getting closer and closer in modern DRAM memory, leading to isolation and interference issues.

\subsubsection{Data Remanence Attacks}\label{sec:security_of_cryptographic_algorithms:attacks:dra}

At a high level, \textbf{data remanence attacks} aim to retrieve data whose residues persist on cryptographic modules after the execution of cryptographic computations. Such residues may be found in both volatile memory (i.e., RAM), non-volatile memory (e.g., flash memory, SSDs, and magnetic hard drives), and even caches and buffers at the processor level. Data remanence attacks may rely on file deletion operations of operating systems (which typically remove entries of deleted files rather than actually overwriting deleted file contents), physical properties of storage media (e.g., RAM) allowing recovery of previous contents, or on applying rapid cooling techniques (e.g., liquid nitrogen) on memory chips to recover data that would have otherwise been erased during a normal shutdown process (in this last case, the attacker is required to be active). Data remanence attacks are often carried out in combination with temperature variations (\Cref{sec:security_of_cryptographic_algorithms:attacks:fia}) \cite{anagnostopoulos_low_2018}. Again, the term data remanence attack denotes a \textbf{category of attacks} encompassing different approaches and techniques for recovering data, e.g., cold boot attacks and forensics recovery. 

\paragraph*{\faSnowflake~Cold Boot Attacks.} This type of data remanence attack relies on the \textbf{physical properties of RAM memory} due to which data persist in RAM memory for a short period (e.g., seconds to minutes) after the system is powered off. In cold boot attacks, the RAM memory is often physically removed from the system and analyzed on a different device. Alternatively, attackers may forcefully and abruptly reboot a system using a lightweight operating system from a removable disk to then dump the contents of the RAM memory \cite{halderman_lest_2009}.

\paragraph*{\faMicroscope~Forensic Recovery.} This type of data remanence attack relies on \textbf{forensics methods} to observe traces of data that should have been deleted and consequently recover the data themselves \cite{weingart_physical_2000}. Forensic recovery attacks exploit residual data traces on rewritable optical media (e.g., CDs, DVDs) and non-volatile memory (i.e., hard drives and tapes through residual magnetic patterns and flash memory or SSDs with wear-leveling).

\remarkBlock{On the practicality of implementation attacks}{Attacks on the implementation of cryptographic algorithms may be more or less practical depending on the context --- which should be reflected by the definition of a suitable threat model (\Cref{sec:cryptographic_security_overview:threat_model}) --- and the system under consideration. In fact, conducting such attacks against complex systems such as laptops and servers is inherently more complex than conducting them against isolated hardware-based cryptographic modules like \gls{se} and \gls{tpm} and embedded systems alike (\Cref{sec:cryptographic_modules}): when the system is small and the hardware to be attacked relatively simple, then these attacks are more likely to succeed. Although it is often the case that software vulnerabilities (e.g., generic malware) are simpler to exploit than vulnerabilities in algorithm implementation, this does not mean that these attacks are impractical. On the contrary, several researchers highlighted the practicality of such attacks --- even for modestly equipped and skilled attackers --- and called for increased awareness among developers and organizations \cite{chiu_sok_2023,dehbaoui_injection_2012,lou_survey_2022,brumley_remote_2003,breier_how_2022}.}
